% SHARD-ID: RC-04D-BC-SYMMETRY-GENERATORS
% SHARD-TITLE: The Phantasm III: BC Stationarity, Arithmetic Symmetry, and the Remaining Generator Cone
% SHARD-SUMMARY: Formulates the finite inverse problem with the BC phase state and Galois/Weil covariance, distinguishing matrix-state extensions from representations of the full BC algebra.
% SHARD-SUMMARY: Computes the exact generator-cone dimensions, proves the extra Weyl-covariance collapse, and gives a symplectic odd-sector counterexample and compatible two-prime couplings.
% SHARD-SUMMARY: Records deterministic numerical checks and the missing representation/intertwiner input needed before any adelic or scattering identification.
% SHARD-KEYWORDS: Bost-Connes, Galois, symplectic, adeles, Weyl, Weil, GKLS, cone, Choi, parity, coupling
% SHARD-NOTES: notes/bc-symmetry-generators.md
% SHARD-SCRIPTS: scripts/bc_symmetry_generators.py

\section{The Phantasm III: Stationarity, Symmetry, and Generator Freedom}
\label{sec:bc-symmetry-generators}

TJO's specification is an unknown Lindbladian with the BC state stationary,
Galois symmetry, and compatible adelic symplectic structure. The no-event
generator and the jumps are both unknown. This section completes a finite
constraint calculation, not the construction of the adelic bond. Detailed
arguments B0--B6 are in \texttt{notes/bc-symmetry-generators.md}; source loci
are in \texttt{notes/extract/bc-symmetry-sources.md}. These arguments have not
had an independent review and therefore carry \texttt{sketched}, including
the statements with complete written derivations. Numerical evidence is
separately tagged. No claim of novelty is made for the general GKLS machinery.

\subsection{The state on the finite phase algebra}

\begin{proposition}[Residue probabilities]
\label{prop:bc-residue-state}
\claimstatus{prop:bc-residue-state}{sketched}
For the state of \cref{def:bc-finite-phase},
\[
 w_{\invtemp}(0)=\pp^{-\invtemp},\qquad
 w_{\invtemp}(x)=\frac{1-\pp^{-\invtemp}}{\pp-1}\quad(x\ne0).
\]
At criticality the phase marginal is uniform at every integer level $N$,
compatible under reduction, and $\BCres{\pp,1}=I/\pp$.
\end{proposition}

This is finite Fourier inversion of the BC KMS formula
\cite{connesmarcolli2004physics}; for $\invtemp>1$ it also follows by averaging
the residue-class Gibbs sums over the Galois action. A single extremal state
is generally not invariant. Uniform phase probabilities fix the diagonal
of a matrix extension, not its off-diagonal entries. The full critical BC
state is not a normal density in the standard representation
(\cref{prop:normal-kms-gibbs}).

\begin{proposition}[Finite representations collapse the BC phases]
\label{prop:bc-no-finite-phase-representation}
\claimstatus{prop:bc-no-finite-phase-representation}{sketched}
Every finite-dimensional unital $*$-representation of the full BC algebra
sends all $e(r)$, $r\in\mathbb Q/\mathbb Z$, to $I$.
Consequently it cannot realise the nontrivial critical finite phase marginal.
\end{proposition}

Indeed, $\bciso{n}^*\bciso{n}=I$ makes $\bciso{n}$ unitary in finite dimension.
The BC relation gives
$I=\bciso{n}\bciso{n}^*=n^{-1}\sum_{k=0}^{n-1}e(k/n)$,
the spectral projection onto the eigenvalue $1$ of $e(1/n)$.
Thus $e(1/n)=I$ for every $n$. This excludes an exact finite representation;
it does not exclude restrictions, CP compressions, or asymptotic models.

\begin{proposition}[Invariant state extensions]
\label{prop:bc-symmetry-state-extension}
\claimstatus{prop:bc-symmetry-state-extension}{sketched}
A unique stationary density of a covariant semigroup is invariant under
its symmetry group. The diagonal extension $\BCres{\pp,\invtemp}$ is
Galois invariant, and is Weil invariant iff $\invtemp=1$.
The same phase marginal has a unique Weil-invariant matrix extension,
when it is positive:
\[
 \BCbond=aI+(r-a)\Par,\qquad
 r=\pp^{-\invtemp},\qquad a=\frac{1-r}{\pp-1}.
\]
Its parity-block eigenvalues are $r$ and $2a-r$, so positivity holds iff
$r\le2/(\pp+1)$. Full Weyl invariance forces $I/\pp$.
\end{proposition}

The Weil-invariant operator space is $\operatorname{span}\{I,\Par\}$:
the symplectic group is transitive on the nonzero Weyl labels and
$\sum_v\Weyl{v}=\pp\Par$. The diagonal fixes the two coefficients.
These coherent extensions are explicit state-extension examples, not an
identification with the full BC bond. Above the transition, covariance can
instead act on a family of stationary extremal states; uniqueness is the
premise of the first assertion.

\subsection{The complete finite feasibility problem}

\begin{theorem}[A linear section of the conditional Choi cone]
\label{thm:bc-covariant-gks-cone}
\claimstatus{thm:bc-covariant-gks-cone}{sketched}
The cone $\Gcone_U(\BCbond)$ consists exactly of complex-linear maps $L$ with
\[
 \begin{gathered}
 L(X^*)=L(X)^*,\quad \Tr L(X)=0,\quad L(\BCbond)=0,\\
 [L,\Ad(U_g)]=0\quad\hbox{for every }g,\qquad Q\Choi_LQ\ge0.
 \end{gathered}
\]
Equivalently, in an orthonormal traceless operator basis $F_i$,
\[
 L(X)=-i[H,X]+\sum_{i,j}A_{ij}
 \left(F_iXF_j^*-\tfrac12\{F_j^*F_i,X\}\right),\qquad A\ge0,\quad H=H^*.
\]
With $\Tr H=0$, covariance requires $H$ invariant and $A$ commuting with
the induced representation on traceless operators. If that representation
is $\bigoplus_\lambda V_\lambda\otimes\mathbb C^{m_\lambda}$, then
$A=\bigoplus_\lambda I_{V_\lambda}\otimes A_\lambda$, $A_\lambda\ge0$;
stationarity supplies further linear equations. Every invariant faithful
$\BCbond$ has an interior feasible generator $X\mapsto\BCbond\Tr X-X$.
\end{theorem}

Necessity follows by differentiating a positive channel Choi matrix on
the complement of the maximally entangled vector. For sufficiency, the
positive matrix $Q\Choi_LQ$ supplies the traceless Kraus part; the remainder
is the Choi matrix of $KX+XK^*$, and trace annihilation fixes its Hermitian
part. This yields the bounded GKLS form recalled in
\cite{siemon2017unbounded}. The reset generator has conditional Choi matrix
$Q(I\otimes\BCbond)Q$, strictly positive on the range of $Q$.

\begin{theorem}[Dimensions at the critical tracial extension]
\label{thm:bc-critical-cone-dimensions}
\claimstatus{thm:bc-critical-cone-dimensions}{sketched}
For odd prime $\pp$, at stationary density $I/\pp$, the real dimensions
of the generator cones' linear spans are
\[
\begin{array}{c|c}
\text{commuting actions}&\text{dimension, including time scale}\\ \hline
\text{Galois torus}&(\pp-1)(\pp+1)^2\\
\text{full linear Weil action}&2\pp-2\\
\text{Weyl translations and Galois torus}&\pp+1\\
\text{Weyl translations and full Weil action}&1.
\end{array}
\]
\end{theorem}

In Weyl coordinates, trace annihilation and stationarity remove the
identity row and column. The first two dimensions count simultaneous
group orbits on ordered pairs of nonzero vectors. The torus acts freely;
the symplectic orbits are $w=av$ ($a\ne0$) and
$\det(v,w)=b$ ($b\ne0$), giving $2(\pp-1)$. Since $-I$ belongs to both
groups, Hermiticity makes the orbit coefficients real. The depolarizing
generator is an interior point of the conditional Choi cone, so positivity
does not reduce these dimensions. Orbit coefficients are not independent
nonnegative rates in the first two rows.

\begin{proposition}[The cost of extra Weyl covariance]
\label{prop:bc-weyl-covariance-collapse}
\claimstatus{prop:bc-weyl-covariance-collapse}{sketched}
All Weyl-covariant generators are
$L(X)=\sum_{v\ne0}c_v(\Weyl{v}X\Weyl{v}^*-X)$, $c_v\ge0$.
Galois covariance makes the rates constant on the two punctured axes and
the $\pp-1$ sets $xy=b\ne0$. In particular $c_v=c_{-v}$, making every
Weyl eigenvalue real and nonpositive. Full symplectic covariance makes
all the rates equal and gives exactly
$L=\gamma(\Dep_{\pp}-\mathrm{id})$, $\gamma\ge0$.
\end{proposition}

This extra commuting symmetry removes the oscillatory zero-like modes in
this model. It is not a consequence of merely requiring a symplectic
representation or an adelic structure.

\subsection{Odd spectral freedom and coupling freedom}

\begin{proposition}[A symplectic odd-sector counterexample]
\label{prop:bc-symplectic-odd-counterexample}
\claimstatus{prop:bc-symplectic-odd-counterexample}{sketched}
Let $a$ be a primitive root modulo $\pp$ and set $R_a(I)=0$,
$R_a(\Weyl{v})=\Weyl{av}$ for $v\ne0$. Then
\[
 L=\Dep_{\pp}-\mathrm{id}+\varepsilon R_a,
 \qquad\varepsilon=\frac1{4\pp^2}
\]
is a CPTP generator with unique stationary state $I/\pp$ and full
Weil, Galois and parity covariance. For $\pp=7$ its odd eigenvalues are
$-1-\varepsilon$ and $-1+\varepsilon/2\pm i\sqrt3\varepsilon/2$, each
with multiplicity $8$. The odd decay rates are unequal.
\end{proposition}

The map $R_a$ is Hermiticity preserving and equivariant, but need not itself
be positive. Its Choi operator norm is bounded by its Frobenius norm
$\sqrt{\pp^2-1}<\pp$, whereas the depolarizer's conditional Choi matrix
is $Q/\pp$; the specified perturbation therefore remains strictly positive
on the range of $Q$. At $\pp=7$, scalar multiplication has eight cycles of
length six, and the odd sector selects the roots whose cube is $-1$.
Thus the counterexample does not fit any input zero positions.

\begin{proposition}[Couplings invisible to the one-prime restrictions]
\label{prop:bc-compatible-prime-coupling}
\claimstatus{prop:bc-compatible-prime-coupling}{sketched}
For distinct primes $p,q$, fixed positive rates $a,b$, and
$0\le c\le\min(a,b)$, let
\[
 \begin{split}
 L_c={}&(a-c)(\Dep_p-\mathrm{id})\otimes\mathrm{id}
 +(b-c)\mathrm{id}\otimes(\Dep_q-\mathrm{id})\\
 &+c(\Dep_p\otimes\Dep_q-\mathrm{id}).
 \end{split}
\]
These CPTP generators have the unique stationary state $I/(pq)$ and full
local Weyl/Weil covariance. Their restrictions on the one-prime observable
algebras are independent of $c$, with rates $a$ and $b$, while the rate on
the joint traceless sector is $a+b-c$.
\end{proposition}

Each summand is a nonnegative multiple of a channel minus the identity.
The four tensor sectors give the eigenvalues $0,-a,-b,-a-b+c$ directly.
This supplies compatible two-prime residue dynamics under CRT with a
variable coupling; it is not an infinite adelic construction.

\begin{numerical}[Finite arithmetic symmetry checks]
\label{num:bc-symmetry-generators}
\claimstatus{num:bc-symmetry-generators}{numerical}
The script \path{scripts/bc_symmetry_generators.py} passes $427$ checks;
its output is \path{outputs/bc_symmetry_generators.txt}. The checks cover
$p=3,5,7$: BC residue probabilities against Hurwitz sums, invariant matrix
extensions, orbit dimensions against independent Burnside counts,
conditional and finite-time Choi positivity, trace preservation,
stationarity and uniqueness, covariance, Egorov, odd-sector spectra, and
two-prime restrictions. The four cone dimensions are respectively
$(32,4,4,1)$, $(144,8,6,1)$, and $(384,12,8,1)$.
\end{numerical}

\subsection{What is still missing}

The finite inverse problem is now a concrete positive-semidefinite cone.
The missing input is a representation or CP comparison map connecting the
full BC algebra, including its prime isometries, to the scattering bond,
with compatible local Weil data. It must select the permitted commuting
symmetries and constrain the generator blocks and the inter-prime couplings.
The $\Gamma_0(N)$ character/scattering computation of
\cref{conj:galois-graded-bond} remains open. Prime-power levels, the real
place, metaplectic compatibility and the critical operator-algebra limit
have not been constructed. A chosen graded transfer spectrum alone does
not identify a physical cMPS norm.

The earlier reset test $-(B\BCbond+\BCbond B^*)\ge0$ is a criterion for a fixed
finite no-event generator $B$ and scalar reset, when its exit flux is
positive. It is not a no-go criterion for the full cone in
\cref{thm:bc-covariant-gks-cone}, where $B$ and all jumps may vary.
